% ch00 前言 (书页 xi–xii / 00_front p013–p014)
\chapter*{前言}
\addcontentsline{toc}{chapter}{前言}
\markboth{前言}{}

在撰写本书第2版时，我力求完成三件事。第一，改正第1版中的全部排印错误。这件事做起来比我预想的要难。虽然我确信我和校对者已改正了所有我们能发现的问题，但这可能仍不够。由于本书无疑还会重印，如果你发现任何错误，请发邮件至 dimer@illinois.edu 告知我。第二，补写第1版中本应包含而我当初放弃撰写的内容，包括格林函数、卢廷格（Luttinger）定理、费米液体短程相互作用的重正化，以及对称性。本着本书是一部物理学而非纯粹技巧或数学的论著之宗旨，这些主题都在首次需要之处就地穿插展开。例如，格林函数一节安排在第7章处理安德森杂质问题之处；为完整起见，卢廷格定理也在同一章的附录中给出。第三，加入反映凝聚态物理中 $\hbar=1$ 研究快速进展的新材料。在这一方面，我依据对学生最有用的判断作了取舍。由于目前没有教科书讲授拓扑绝缘体和莫特绝缘体的教学内容（尽管有一些优秀的综述文章），我选择聚焦于这些主题。在撰写拓扑绝缘体部分时，我力图采用所需定义和新概念最少的表述方式，因为这些体系的物理本质上就是简单的。关于莫特问题，我给出的内容自认为是毫无争议的，但尚未在任何单一文献中被系统地写下来。第16章从刚性能带图像成立的能带绝缘体出发，进而论证莫特问题的物理之所以特立独行，正是因为刚性能带图像在其中完全失效。关于刚性能带模型的专著层出不穷，却没有一本书讨论刚性能带图像在强关联电子问题中的失效。铜氧化物问题正是在这一语境下讨论的。我还曾计划写一章量子计算，并把第14章的讨论扩展到玻色–哈伯德模型，但这些内容会把篇幅推到远超600页，使本书过于臃肿。而且依我之见，这些主题并不特别适合一学期的核心研究生课程，而更适合更专门的课程。也许几年后我的想法会有所改变。

在本书手稿的最终编辑过程中，我得益于许多人的帮助。Babak Seradjeh、Juan Jottar 和 Taylor Hughes 对拓扑绝缘体一节提出了宝贵的批评。李威成（Wei-Cheng Lee）、Mohammad Edalati 和 Taylor Hughes 还阅读了莫特一章，找出了若干排印错误和不准确之处。我也感谢 Mohammad 阅读并订正了对称性一章，感谢 Robert Leigh 以其一贯冷静而犀利的风格对强耦合物理和对称性提出的意见。Wade deGottardi 对玻色化一章提出了大量建议。除了世界各地多位学生来信指出他们发现的排印错误之外，我要特别感谢魏伟（Wei Han），他发现了第1版两幅插图中的两处关键错误，也多谢 Taylor Hughes 重绘了这些图。校对的重任落在了我的研究小组以及厄巴纳 ICMT 小组其他成员的肩上——他们的手臂至今仍在为我施加给他们的非绝热扭曲而酸痛。他们包括李威成（Wei-Cheng Lee）、Mohammad Edalati、Seungmin Hong、Wade deGottardi、Rodrigo Garrido 和 Kiaran Dave。此外，在校样阶段，Kiaran Dave、Ka Wai Lo 和郑辉辉（Huihuo Zheng）通读了全书手稿，并以消除一切排印错误为己任，一丝不苟地加以订正。我还要感谢 Matthew Feickert 将 \LaTeX{} 源文件转换为剑桥出版社的格式，并在此过程中发现了若干排印错误；感谢剑桥出版社的员工 Mike Nugent、Simon Capelin、Claire Poole、Abigail Jones 和 Frances Nex 对本项目的热忱投入。一些早年影响对本书的诞生不可或缺：我的高中英语老师 Duane Kusler 鼓励我写作，我的双胞胎妹妹 Andi 和 Lyndi 则让我学到了许多数学小技巧。最后，谨向始终支持我、给我以宁静力量的家人致以由衷的谢意。
